%% Related work (MedIA revision). Every factual statement about a cited paper is
%% traced to its source in sections/related_claims.md.
\section{Related work}
\label{sec:related}

\subsection{Site and center effects}
\label{sec:related-site}

Performance loss at institutions not seen in training is documented across
medical imaging. \citet{zech2018variable} trained pneumonia classifiers on chest
radiographs from three hospital systems. Internal performance exceeded external
performance in three of five comparisons, and for two of the systems the
networks identified the system a radiograph came from with more than 99.9\%
accuracy. In a systematic review of 86 externally validated radiology
algorithms, 81\% performed worse on external data and 24\% lost at least 0.10
on the unit scale \citep{yu2022external}. In histopathology,
\citet{howard2021site} showed that the submitting site of TCGA slides can be
identified by deep networks even after common color normalization and
augmentation. These site signatures biased the accuracy of survival, mutation
and stage prediction. The authors proposed splits in which no site contributes
to both training and validation. \citet{stacke2019closer} proposed a
model-specific distance between domains, computed in the learned
representation. \citet{stacke2021measuring} developed it into the
\emph{representation shift} and showed that it correlates highly with the drop
in performance across many types of domain shift. These studies show that site
effects exist and can be detected, but not how much of the spread in reported
performance they account for relative to the training method. Our variance
decomposition estimates this.

Mitotic-figure detection has been studied directly under domain shift. The
MIDOG~2021 challenge trained on 200 breast cancer cases from four scanners (one
of them without labels) and
tested on 100 cases from four scanners, two of which were not seen in training.
The winning method reached an F1 score of 0.748 \citep{aubreville2023midog}.
MIDOG~2022 extended the shift across tumor types, laboratories and species
(as its title states) and evaluated nine methods on ten independent domains. Domain characteristics
absent from the training data (a new species, a new cell morphology and a new
scanner) led to small but significant drops in performance
\citep{aubreville2024midog2022}. The method that won both editions combined a
two-stage detector with domain-generalization measures
\citep{jahanifar2024mitosis}. MIDOG++ \citep{aubreville2023midogpp} released
503 images of seven tumor types and found that leave-one-domain-out training
generalized considerably better than training on a single domain. The
multi-scanner canine cutaneous squamous cell carcinoma dataset
\citep{wilm2023canine} digitized the same slides with five scanners, which
isolates the acquisition shift. We use these cohorts, but our question is not
the challenges' question of which detector is best. We ask how much of a
comparison between methods is decided by the domains themselves.

\subsection{Domain generalization benchmarks and protocols}
\label{sec:related-dg}

Leave-one-domain-out evaluation with repeated runs is established practice in
domain generalization (DG); \citet{zhou2023domain} survey the field. DomainBed
\citep{gulrajani2021domainbed} evaluated fourteen algorithms on seven
multi-domain datasets. It trained on the remaining domains for every choice of held-out test
domain, and it repeated the whole study three times, redrawing hyperparameters,
weight initializations and data splits. The authors argued that a DG algorithm
without a model-selection rule is incomplete. They compared three rules:
validation data pooled from the training domains, leave-one-domain-out
cross-validation over the training domains, and an oracle validation set drawn
from the test domain. With training-domain or oracle selection, no algorithm
improved on the average accuracy of empirical risk minimization (ERM) by more
than one point, and the authors found no improvement over ERM statistically
significant. Under leave-one-domain-out selection, however, CORAL led ERM by
1.8 points \citep[Table~3]{gulrajani2021domainbed}.
Training-domain validation also outperformed leave-one-domain-out
cross-validation across multiple datasets and algorithms.

WILDS \citep{koh2021wilds} made Camelyon17 \citep{bandi2019camelyon17} a DG
benchmark with three training hospitals, one out-of-distribution (OOD)
validation hospital and one test hospital. The test hospital was chosen because
its patches were the most visually distinctive. ERM test accuracy had a
standard deviation of 6.4 points over ten seeds, and CORAL, IRM and GroupDRO
performed comparably to or worse than ERM. WILDS therefore asks for ten seeds
per Camelyon17 submission and warns that so few test domains invite overfitting
to them. WILDS~2.0 \citep{sagawa2022extending} added unlabeled patches from the
training, validation and test hospitals. On Camelyon17, ERM reached 70.8\% OOD
accuracy without data augmentation and 82.0\% with it. SwAV pre-training on
unlabeled target-hospital patches reached 91.4\% (standard deviation 2.0) and
Noisy Student reached 86.7\%. CORAL, DANN, Pseudo-Label and FixMatch all fell
below augmented ERM \citep[Table~2]{sagawa2022extending}. These figures come
from one held-out hospital. The best arm also sees unlabeled images from that
hospital, so it is unsupervised adaptation rather than DG.
\citet{wiles2022fine} distinguished spurious correlation, low-data drift and
unseen data shift, and trained more than 85,000 models covering 19 methods on
six datasets, Camelyon17 among them. Pre-training and augmentation often helped
substantially, but the best methods differed across datasets and shifts. On
Camelyon17, color jitter hurt performance, and pre-training was not the best
choice under unseen data shift.

We adopt leave-one-domain-out evaluation and training-domain model selection
from this literature and claim neither as new. Our contribution is diagnostic.
We decompose performance into between-domain and between-run variance with
interval estimates, and we analyze how stable method rankings are across
held-out domains and cohorts. We use a factorial cohort design that crosses task
(tumor versus mitotic figure) with shift type (institutional versus
acquisition-only). Ranking agreement can then be compared between cohorts that
share a task and cohorts that share a shift type.

\subsection{Stain normalization and augmentation}
\label{sec:related-stain}

Differences in staining between laboratories degrade networks applied to
slides from an unseen laboratory \citep{tellez2019stain}.
\citet{reinhard2001color} transferred color by matching per-channel means and
standard deviations in the decorrelated $l\alpha\beta$ space. Color
deconvolution \citep{ruifrok2001deconvolution} separates up to three stains in
RGB images using stain-specific absorption, and it is the basis of HED-space
augmentation \citep{tellez2019stain}. \citet{macenko2009} estimate the stain
vectors of each slide from its optical densities. Pixels are projected onto the
plane of the two leading singular vectors, the vectors are taken at robust
angular extremes, and per-stain intensities are rescaled to a common
pseudo-maximum. \citet{vahadane2016structure} factorize images into sparse,
non-negative stain density maps and recombine them with a target image's stain
color basis. StainGAN \citep{shaban2019staingan} learns the mapping with a
CycleGAN-style model, so no single reference slide has to be chosen. Instead of normalizing,
\citet{tellez2018whole} perturbed the hematoxylin and eosin channels directly
and, together with ensembling, obtained a stain-invariant mitosis detector.
Across four tasks and nine laboratories, \citet{tellez2019stain} found HSV or
HED color augmentation essential. The top-ranked configuration used no
normalization, although normalization made classifiers less sensitive to the
choice of augmentation. Normalization also leaves site information in place
\citep{howard2021site}. We therefore evaluate per-slide Macenko normalization,
a per-tile variant, H\&E jitter and their combination as separate arms.

\subsection{Self-supervision and pathology foundation models}
\label{sec:related-fm}

\citet{ciga2022ssl} pre-trained SimCLR on 57 unlabeled histopathology datasets.
Linear classifiers on these features beat ImageNet features by more than 28\%
in F1 on average. \citet{kang2023benchmarking} compared MoCo~v2,
SwAV, Barlow Twins and DINO pre-trained on 19 million patches from 20,994 TCGA
slides, and domain-aligned pre-training outperformed ImageNet pre-training in
both linear and fine-tuning evaluation. Foundation models extend this further:
\begin{itemize}
  \item CTransPath, a convolutional network combined with a multi-scale Swin
    Transformer \citep{wang2022ctranspath};
  \item Phikon, an iBOT ViT-Base trained on more than 40 million images from 16
    cancer types \citep{filiot2023phikon};
  \item UNI, trained on more than 100 million images from over 100,000
    whole-slide images (WSIs) covering 20 tissue types \citep{chen2024uni};
  \item CONCH, trained on more than 1.17 million image--caption pairs
    \citep{lu2024conch};
  \item Virchow, a 632-million-parameter ViT trained with DINOv2 on about
    1.5 million WSIs \citep{vorontsov2024virchow};
  \item Prov-GigaPath, trained on 1.3 billion tiles from 171,189 slides from a
    network of 28 cancer centers \citep{xu2024gigapath}.
\end{itemize}

Scale does not by itself remove center effects. \citet{komen2024batch} found
hospital signatures in foundation-model embeddings that dominated feature-space
distances and were not removed by stain normalization. All ten public
foundation models evaluated by \citet{dejong2025unrobust} encoded the medical
center strongly. Only one had a robustness index above one, meaning that
biological features outweighed center features, and then only slightly.
The PathoROB benchmark found robustness deficits in all 20 models it evaluated
\citep{komen2026robust}. More robust models, vision--language alignment and
post-hoc robustification reduced, but did not eliminate, the risk of the
resulting downstream errors. Our analysis is complementary. We measure how the
downstream performance of frozen and fine-tuned pathology encoders
\citep{kang2023benchmarking} varies across held-out domains, under the same
protocol and diagnostics as every other arm.

\subsection{Test-time adaptation}
\label{sec:related-tta}

AdaBN adapts a trained network without labels by replacing its
batch-normalization (BN) statistics with target-domain statistics, and it adds
no parameters \citep{li2017adabn}. \citet{schneider2020bn} showed that
re-estimating BN statistics on corrupted images improved ImageNet-C robustness
for 25 architectures, and argued that adapted results should be reported in OOD
evaluations. Prediction-time BN improved accuracy and calibration under
covariate shift, but its results were mixed when combined with pre-training and
weaker under more natural shifts \citep{nado2020bn}. TENT \citep{wang2021tent}
minimizes prediction entropy online, updating normalization statistics and
channel-wise affine parameters on each batch. Because these methods estimate
statistics from test batches, which depend on batch composition, our protocol
shuffles the test stream with a fixed seed before any adaptation.

\subsection{Variance, rankings and statistical reporting}
\label{sec:related-stats}

Data sampling, parameter initialization and hyperparameter choice each have a
marked effect on benchmark results \citep{bouthillier2021variance}, and
test-set scores alone do not settle which model is best \citep{dodge2019show}.
In biomedical image analysis challenges, an algorithm's rank is generally not
robust to the test data, the ranking scheme or the annotators
\citep{maierhein2018rankings}. challengeR provides bootstrap analyses of ranking
stability \citep{wiesenfarth2021challenger}. \citet{varoquaux2018cv} showed
that cross-validation error bars are large (about $\pm 10\%$ with 100 samples)
and that the standard error across folds strongly underestimates them. In two
Kaggle medical-imaging challenges with test sets below 1,000 cases, the
difference between the public and private leaderboards was larger than the gap
between the winner and the top-10\% entry \citep{varoquaux2022failures}. Any variance estimator computed only from the
results of a cross-validation experiment is biased, and ignoring
training-set variability can make an algorithm look significantly better than
it is \citep{nadeau2003inference}. Cross-validation estimates the average error
of models fit to other training sets from the same population (proved for
least-squares linear models), and the correlation between folds makes
the usual intervals too narrow \citep{bates2023cv}. Leave-one-domain-out
designs are an extreme case. There are few folds ($K=3$ to $7$ here), the folds
share training domains, and the held-out domain is both a nuisance factor and
the quantity of interest. We therefore report between-domain and between-run
variance components with intervals and run a fold-dependence sensitivity
analysis on paired contrasts. We also give bootstrap intervals for method ranks
(over seeds and over domains) and for Kendall's $\tau$ between cohort rankings.
